% Introducción: cómo leer la solución, dónde están los retos y reparto del tiempo
\section*{Introducción}

\begin{itemize}
\item \textbf{\emph{Qué es}}: la solución \emph{de máximos} del examen de ejemplo del dictamen sobre España (10 preguntas, una por bloque de la guía, 3 horas y 30 minutos orientativas). Aplica la guía de principio a fin y remite a sus apartados («guía, I.4» es el apartado~I.4, \emph{Demanda}). Todas las cifras son del cuadernillo (cuadros 1 a 9 y 2~bis, gráficos 1 a 14, datos oficiales a 26 de octubre de 2026); lo que viene de fuera se dice y se cita.
\item \textbf{\emph{Cómo se lee}}: cada pregunta tiene 3 piezas, que se distinguen a simple vista:
  \begin{itemize}
  \item el \textbf{enunciado}, en un recuadro verde, con el \emph{peso} de la pregunta, su tiempo y la extensión del texto;
  \item el \textbf{texto de examen}, con una \textcolor{tceemorado}{\textbf{raya morada}} a la izquierda: lo que escribiría un opositor excelente, con rótulos para cada subpregunta, frase de diagnóstico al empezar y cifras citadas por su cuadro o gráfico;
  \item los \textbf{cuadros complementarios}, en gris con una \textcolor{tceedorado}{\textbf{raya dorada}}: por qué se dice cada cosa, qué apartado de la guía se aplica, los cálculos hechos con el cuadernillo, gráficos que el examen no trae, errores típicos, otras formas de redactarlo y qué añadir con más tiempo. \uline{No forman parte de la respuesta}: son la clase que hay detrás.
  \end{itemize}
\item \textbf{\emph{Un techo, no un objetivo}}: el texto de examen es más largo de lo que cabe a mano en el tiempo de cada pregunta. Está escrito para que se vea todo lo que se \emph{puede} decir; en el examen se conserva la frase de diagnóstico de cada párrafo y 1 o 2 cifras por idea (guía, D.5), y cada pregunta termina con un cuadro que dice qué recortar.
\end{itemize}

\paragraph{Dónde están los retos y cuánto pesa cada pregunta}
\begin{itemize}
\item La guía pide repartir el tiempo según los retos reales de \emph{esa} economía, no a partes iguales (guía, introducción y D.3). Leído el cuadernillo, el diagnóstico de España cabe en una frase: \uline{crece más que la zona del euro, con superávit exterior, banca rentable y deuda privada en mínimos, pero con la productividad estancada, un paro que dobla al europeo, una deuda pública en torno al 100~\% del PIB y un nuevo choque de precios en 2026}.
\item De ahí los pesos: \textbf{alto} en el perfil cíclico (ordena todo y recoge un ciclo completo con 2 recesiones), el mercado de trabajo y la productividad, la inflación, las cuentas públicas y las recomendaciones; \textbf{medio} en la inversión y el crédito, el sector exterior y la política monetaria (hay mucho que contar, pero hoy no son un problema o no se deciden en España); \textbf{bajo} en la competitividad no precio y la banca (exposición pequeña a los aranceles y banca sólida).
\item La extensión sigue a los pesos: \uline{alto~$\gg$~medio~$>$~bajo}. La tabla~\ref{tab:reparto} aplica la regla de la guía (lectura y repaso fijos; el resto, $t=W/(n+0{,}25)$ y $1{,}25\,t$ para el perfil cíclico) y desplaza minutos de las preguntas de peso bajo a las de peso alto.
\end{itemize}

\begin{table}[H]
\centering\footnotesize
\renewcommand{\arraystretch}{1.12}
\caption{Reparto orientativo del tiempo y extensión del texto de examen según el peso de cada pregunta}\label{tab:reparto}
\begin{tabular}{>{\raggedright\arraybackslash}p{4.1cm}>{\centering\arraybackslash}p{1.1cm}>{\centering\arraybackslash}p{1.25cm}>{\centering\arraybackslash}p{1.5cm}>{\raggedright\arraybackslash}p{6.4cm}}
\toprule
\rowcolor{tceeazulpalido}
\textbf{Pregunta (bloque de la guía)} & \textbf{Peso} & \textbf{Minutos} & \textbf{Palabras} & \textbf{Por qué ese peso} \\
\midrule
Lectura del cuadernillo y plan & — & \tiempo{15} & — & Fijo: fases marcadas, cálculos hechos y gráficos emparejados con su pregunta (cuadro complementario~0.1) \\
\rowcolor{tceesalvia}
1. Perfil cíclico (I) & alto & \tiempo{23} & 2.100% & Ordena todo el examen; un ciclo completo con 2 recesiones; el «porqué» del PIB por habitante \\
\rowcolor{tceesalvia}
2. Trabajo y productividad (II.1) & alto & \tiempo{21} & 1.120% & Paro estructural y productividad estancada: los 2 grandes retos de España \\
3. Consumo, inversión y crédito (II.2) & medio & \tiempo{17} & 1.408% & Buena historia (desapalancamiento, inversión más sana que en 2007), pero no es un problema; el foco, la vivienda \\
\rowcolor{tceesalvia}
4. Inflación y competitividad precio (III.1-2) & alto & \tiempo{21} & 1.579% & El problema coyuntural de 2026: 5,0~\% frente a 3,8~\%, riesgo de segunda ronda y CLU por encima de la zona \\
5. No precio y aranceles (III.3) & bajo & \tiempo{13} & 508% & Exposición directa a EE.~UU. pequeña y cuotas estables; el reto de fondo (I+D) va a la 10 \\
6. Sector exterior (IV) & medio & \tiempo{17} & 760% & La gran corrección está hecha y la PIIN mejora; quedan la energía y una PIIN aún bajo el $-35$~\% \\
7. Política monetaria (V.1) & medio & \tiempo{17} & 709% & El tono no lo decide España; sí la transmisión rápida y un tipo real más bajo que en la zona \\
8. Banca y macroprudencial (V.2) & bajo & \tiempo{13} & 432% & Banca rentable y solvente, deuda privada en mínimos; vigilar la vivienda \\
\rowcolor{tceesalvia}
9. Cuentas públicas y fiscal (V.3-4) & alto & \tiempo{21} & 1.137% & Deuda en torno al 100~\%, déficit estructural del 3~\% con brecha positiva, envejecimiento y reglas nuevas \\
\rowcolor{tceesalvia}
10. Desequilibrios y recomendaciones (VI) & alto & \tiempo{21} & 1.081% & La síntesis: donde se juntan productividad, paro, deuda y demografía \\
Repaso & — & \tiempo{11} & — & Releer, completar la frase de diagnóstico de la pregunta 1 y cerrar lo abierto \\
\midrule
\textbf{Total} & & \textbf{3 h 30 min} & \textbf{8.782}% & \\
\bottomrule
\end{tabular}
\fuente{Elaboración propia con la regla de la guía (D.3): $W=210-15-11=184$ minutos; con $n=10$, $t=184/10{,}25\approx18$ y $1{,}25\,t\approx22$. Dentro de esa media, cada pregunta de peso alto recibe 21 minutos (23 el perfil cíclico), las de peso medio 17 y las de peso bajo 13 (suman 184). Palabras: las del texto de examen de esta solución (sin enunciados ni cuadros complementarios). En verde, las preguntas de peso alto.}
\end{table}

\begin{complemento}{0.1 · Los 15 minutos de lectura: qué se marca en el cuadernillo}
\que{Qué se hace} (guía, D.4): leer las 10 preguntas enteras subrayando lo que pide cada subpregunta; mirar las unidades; dividir el ciclo con el PIB y la brecha (cuadro~1); marcar en verde y rojo fortalezas y debilidades; emparejar cada cuadro y gráfico con su pregunta, y hacer al principio los cálculos que se van a usar.

\que{Emparejar cada pieza con su pregunta.} Ningún gráfico está de adorno (guía, D.4): la tabla~\ref{tab:emparejar}, al final de este cuadro, lo hace para este examen.


\que{Los cálculos que conviene hacer al principio} (porque se usan en varias preguntas y así no hay contradicciones): el crecimiento del potencial que implica la brecha; la productividad y los CLU de 2025 ($4{,}4-0{,}0=4{,}4$); el margen ($3{,}0-4{,}4=-1{,}4$); el tipo real ($2{,}50-5{,}0=-2{,}5$); el esfuerzo fiscal ($\Delta b^{pe}=+0{,}3$ en 2025) y el saldo primario que estabiliza la deuda ($\approx-3{,}1~\%$ del PIB). Todos están desarrollados en los cuadros de cada pregunta y en el fichero \texttt{calculos.txt}.

\que{Unidades que engañan en este cuadernillo}: en el tipo de cambio efectivo real (C6, G8) un aumento es una \emph{apreciación}; la deuda externa neta va con signo positivo en C7 y negativo en G11; en la cuenta financiera, $+$ es salida neta de capital; la brecha, el potencial y los saldos ajustados son de AMECO (primavera de 2026), anterior a la última revisión del INE (PIB de 2025: 2,6~\% en C1 y 2,8~\% en AMECO).

\begin{tablacaja}
\centering\scriptsize
\renewcommand{\arraystretch}{1.08}
\caption{Qué cuadros y gráficos alimentan cada pregunta}\label{tab:emparejar}
\begin{tabular}{>{\raggedright\arraybackslash}p{3.3cm}>{\raggedright\arraybackslash}p{4.6cm}>{\raggedright\arraybackslash}p{6.3cm}}
\toprule
\rowcolor{tceeazulpalido}
\textbf{Pregunta} & \textbf{Piezas del cuadernillo} & \textbf{Para qué} \\
\midrule
1. Perfil cíclico & C1, C2, C2 bis, G1, G2 & Fases, contribuciones, oferta, 2026 trimestral, PIB por habitante \\
2. Trabajo & C3, C2 bis, G3 & Paro, Beveridge, temporalidad, productividad, CLU \\
3. Capital y crédito & C1 (FBCF), C4, G4, G5 & Ahorro, inversión por activos, deuda privada, vivienda \\
4. Inflación & C5, C3 (CLU), C6, G6, G7, G8 & Brote de 2021-2023, choque de 2026, segunda ronda, TCER \\
5. No precio & C6, G9 & Cuotas, I+D, alta tecnología, exportaciones a EE.~UU. \\
6. Sector exterior & C7, C4 (ahorro e inversión), G10, G11 & Subcuentas, financiación, PIIN y efecto valoración \\
7. Monetaria & C8, C4 (tipos), G12, G6 & Tono (tipo real), canales, reacción de 2026 \\
8. Banca & C8, C4, G4 & Semáforo de ratios, vivienda, deuda privada \\
9. Fiscal & C9, C1 (brecha), G13, G14 & Saldos, esfuerzo, $i-\gamma^N$, reglas \\
10. Recomendaciones & todo (matriz de fortalezas y debilidades) & Síntesis y riesgos de las previsiones \\
\bottomrule
\end{tabular}
\fuente{Cuadernillo del examen de ejemplo. C: cuadro; G: gráfico.}
\end{tablacaja}
\end{complemento}
